% Демонстрация 3: «Калькулятор расстояния до плоскости и угла» (Лекция 10)
\section*{Демонстрация 3. «Расстояние до плоскости и угол»\\ \normalsize\textmd{Файл \texttt{demos/03-plane-distance.html} \quad(к Лекции~10)}}

\subsection*{Какие факты лекции иллюстрирует}

Демонстрация представляет собой интерактивную 3D-сцену, в которой
преподаватель и студенты в реальном времени управляют плоскостью,
точкой и прямой, наблюдая за пересчётом всех геометрических величин.
Она наглядно связывает аналитические формулы Лекции~10 с их
геометрическим содержанием.

\begin{itemize}
  \item \textbf{Нормаль как коэффициенты уравнения.}
    Ползунки $A$, $B$, $C$ задают координаты нормального вектора
    $\bar{n}(A,B,C)$, а синяя стрелка в сцене сразу разворачивается в
    соответствующую сторону. Перетаскивая нормаль, студенты видят, что
    коэффициенты общего уравнения
    \[
      Ax + By + Cz + D = 0
    \]
    --- это и есть проекции нормали к плоскости, а ползунок $D$ лишь
    сдвигает плоскость вдоль $\bar{n}$, не меняя её ориентации.

  \item \textbf{Формула расстояния и её геометрический смысл.}
    При любом положении точки $\bar{M}_0(x_0,y_0,z_0)$ панель выводит
    \[
      d = \frac{|Ax_0 + By_0 + Cz_0 + D|}{\sqrt{A^2 + B^2 + C^2}},
    \]
    а в сцене строится зелёный перпендикуляр из $\bar{M}_0$ в красную
    точку проекции $\bar{M}_0'$. Длина этого отрезка совпадает с
    числовым значением $d$. Тем самым видно, что расстояние --- это
    \emph{модуль проекции} вектора $\overrightarrow{M_1M_0}$ на
    единичную нормаль $\bar{n}/|\bar{n}|$ (ровно так оно и выводилось в
    доказательстве теоремы о расстоянии).

  \item \textbf{Знак выражения и полупространства.}
    Числитель формулы стоит под модулем, но \emph{внутреннее} выражение
    $Ax_0 + By_0 + Cz_0 + D$ имеет знак: он положителен по ту сторону
    плоскости, куда смотрит нормаль, и отрицателен --- по другую.
    Передвигая точку сквозь плоскость, студенты наблюдают, как $d$
    убывает до нуля и снова растёт, в то время как точка переходит из
    одного полупространства в другое.

  \item \textbf{Угол между прямой и плоскостью.}
    Оранжевый вектор $\bar{s}(s_x,s_y,s_z)$ задаёт прямую, проходящую
    через $\bar{M}_0$. Панель пересчитывает угол $\varphi$ по формуле
    \[
      \sin\varphi
      = \frac{|(\bar{n}, \bar{s})|}{|\bar{n}|\,|\bar{s}|}
      = \frac{|As_x + Bs_y + Cs_z|}
             {\sqrt{A^2+B^2+C^2}\,\cdot\,\sqrt{s_x^2+s_y^2+s_z^2}}.
    \]
    Здесь наглядна ключевая тонкость лекции: скалярное произведение
    даёт \emph{косинус} угла между $\bar{n}$ и $\bar{s}$, но это угол с
    \emph{нормалью}, а сам угол прямой с плоскостью дополняет его до
    $\pi/2$ --- поэтому в формуле появляется синус.
\end{itemize}

\begin{remarkIT}{Где это работает за пределами доски}{demo03_applications}
  Расстояние от точки до плоскости и проекция --- одни из самых
  «горячих» операций в прикладных 3D-вычислениях:
  \begin{itemize}
    \item \textbf{Коллизии в 3D-движках.} Проверка, не прошёл ли объект
      сквозь стену или пол, сводится к вычислению $d$ и контролю знака
      выражения $Ax_0+By_0+Cz_0+D$: смена знака означает пересечение
      плоскости.
    \item \textbf{Clipping и отсечение по плоскости.} Графический
      конвейер отбрасывает геометрию вне пирамиды видимости камеры,
      сравнивая каждую вершину с шестью плоскостями отсечения по тому
      же знаку.
    \item \textbf{Лидар и навигация роботов.} Выделение опорной
      плоскости (пол, дорога) из облака точек и привязка к ней
      измерений --- это массовое вычисление расстояний $d$; угол
      $\varphi$ между лучом сенсора и плоскостью определяет качество
      измерения.
  \end{itemize}
  Все эти алгоритмы --- буквально формула расстояния $d$ и формула угла
  $\varphi$ из этой демонстрации, исполняемые миллионы раз в секунду.
\end{remarkIT}

\subsection*{Примерный сценарий использования на занятии}

Демонстрацию удобно включать сразу после доказательства теоремы о
расстоянии от точки до плоскости и перед разбором примеров. Ниже ---
один из возможных сценариев.

\begin{enumerate}
  \item \textbf{Запуск и узнавание формулы.} Откройте демонстрацию на
    значениях по умолчанию. Покажите панель «Формулы»: это в точности
    формула $d$ и $\sin\varphi$, только что выведенные на доске. Задача
    --- увидеть в живой сцене все объекты доказательства: нормаль,
    точку, перпендикуляр, проекцию.

  \item \textbf{Нормаль и коэффициенты.} Двигайте ползунки $A$, $B$,
    $C$ и обратите внимание студентов на синюю стрелку: плоскость
    разворачивается вслед за нормалью. Вопрос аудитории: \emph{«что
    задаёт тройка $(A,B,C)$ --- точку плоскости или её наклон?»}

  \item \textbf{Сдвиг плоскости.} Меняйте только $D$. Плоскость
    скользит параллельно самой себе. Подчеркните: ориентация (а значит,
    и нормаль) не зависит от $D$ --- меняется лишь свободный член.

  \item \textbf{Измерение расстояния.} Зафиксируйте плоскость и двигайте
    точку $\bar{M}_0$. Сопоставьте длину зелёного перпендикуляра с
    числом $d$ на панели. Вопрос: \emph{«что будет с $d$, если точка
    окажется на плоскости?»} --- подведите аудиторию к ответу $d=0$.

  \item \textbf{Момент удивления: смена знака.} Медленно проведите точку
    \emph{сквозь} плоскость. Студенты видят, как $d \to 0$, перпендикуляр
    стягивается в точку, а затем расстояние снова растёт --- но точка уже
    в \emph{другом} полупространстве. Обсудите: модуль в формуле «прячет»
    знак, а само выражение $Ax_0+By_0+Cz_0+D$ его сохраняет и отвечает на
    вопрос «по какую сторону». Это и есть основа теста коллизий.

  \item \textbf{Связь с выводом формулы.} Вернитесь к доказательству:
    перпендикуляр в сцене --- это и есть проекция $\overrightarrow{M_1M_0}$
    на нормаль, а $d$ --- её длина. Демонстрация делает абстрактную
    выкладку осязаемой: формула не «придумана», а измерена.

  \item \textbf{Прямая и угол.} Включите вектор $\bar{s}$ и двигайте его
    ползунки. Следите за углом $\varphi$ на панели. Вопрос: \emph{«когда
    угол $\varphi = 90^\circ$?»} --- подведите к тому, что это
    происходит, когда $\bar{s} \parallel \bar{n}$, то есть прямая
    перпендикулярна плоскости. Затем добейтесь $\varphi = 0$: прямая
    легла в плоскость, $\bar{s} \perp \bar{n}$, и скалярное произведение
    обнулилось.

  \item \textbf{Переход к примерам.} Перед разбором задач на проекцию и
    угол попросите студентов \emph{предсказать} знак выражения и
    приблизительную величину угла «на глаз» по сцене, а затем проверить
    расчётом. Это закрепляет связь «формула $\leftrightarrow$ картинка»
    и готовит к самостоятельному решению.
\end{enumerate}
